\section{Evaluation protocol}
\label{sec:evaluation}

\subsection{Partitioning}

We use leave-last-one-out per user. Ordering each user's interactions by the
temporal proxy of Section~\ref{sec:data}, the final interaction forms the test
instance, the penultimate one the validation instance, and everything earlier
is training data.

Users with fewer than $k_{\min} + 2$ interactions cannot yield both held-out
instances without starving their own training history; these users contribute
entirely to the training partition and are not evaluated. Reporting metrics
over users whose training history is empty would measure the popularity prior
rather than personalisation.

\subsection{Candidate set}

For each evaluated user we rank the \emph{entire} catalogue minus that user's
training history:
\begin{equation}
  C_u = \Iset \setminus \Iset_u^{+,\text{train}} .
\end{equation}

Excluding the history is not a detail. Without it, a model is rewarded for
re-recommending already-consumed titles, which inflates every metric and
corresponds to a product behaviour nobody wants. A test constructs a model
that returns only the user's history and asserts it scores exactly zero.

We rank the full catalogue rather than sampling $100$ negatives per user.
Sampled-negative evaluation is common but is known to produce rankings of
methods that disagree with the full-catalogue result~\cite{krichene2020}.

\subsection{Metrics}

With exactly one relevant item $i^{*}_u$ per evaluated user and
$\mathrm{rank}_u \in \{1, \dots, |C_u|\}$ its position in the predicted list,
the cut-off metrics over the evaluated user set $\Uset_{\text{eval}}$ are:

\begin{align}
  \mathrm{NDCG@}K &= \frac{1}{|\Uset_{\text{eval}}|} \sum_{u}
    \frac{\mathbb{1}[\mathrm{rank}_u \leq K]}{\log_2(\mathrm{rank}_u + 1)} ,
    \label{eq:ndcg} \\[0.3em]
  \mathrm{Recall@}K &= \frac{1}{|\Uset_{\text{eval}}|} \sum_{u}
    \mathbb{1}[\mathrm{rank}_u \leq K] , \\[0.3em]
  \mathrm{MRR@}K &= \frac{1}{|\Uset_{\text{eval}}|} \sum_{u}
    \frac{\mathbb{1}[\mathrm{rank}_u \leq K]}{\mathrm{rank}_u} , \\[0.3em]
  \mathrm{Coverage@}K &= \frac{1}{|\Iset|}
    \Big\lvert \bigcup_{u \in \Uset_{\text{eval}}} \mathrm{top}_K(u) \Big\rvert .
\end{align}

Equation~\ref{eq:ndcg} uses the general NDCG definition specialised to a
single relevant item, where the ideal DCG equals $1$ and the normalisation
therefore disappears. Under this protocol Recall@$K$ coincides numerically
with hit rate; both are reported so the table stays comparable with
multi-positive projects.

Coverage is reported because it is the business reason to leave a popularity
list, but it must be read alongside accuracy: the random baseline attains
coverage $1.0$ while being useless.

\subsection{Model selection}

All hyperparameters --- for the baselines as well as for BPR --- are selected
by grid search on the \emph{validation} partition. The test partition is read
once, to produce Table~\ref{tab:results}.

This is stated explicitly because the alternative is both common and
invisible in a results table. Selecting on test turns the reported number
into an optimistically biased estimate whose gap from reality grows with the
size of the grid.

\subsection{Training/serving parity}

Offline metrics are worthless if production serves a different ranking. The
serving path precomputes top-$N$ lists into DynamoDB, and
\texttt{tests/test\_features\_parity.py} asserts that the stored list for each
user is element-wise identical to the list produced by the same fitted model
during evaluation, that no served list intersects the user's history, and
that the payload is JSON-serialisable --- NumPy scalars are a frequent and
late-surfacing cause of rejected writes.
